The results of the article can be followed with a calculator or with the workbook (sheet \texttt{EQ\_P1}). Each part gives the idea, the formula and a worked example.

\subsection{1. the water balance of a soil column}
Over a year, water in (rain $P$ and irrigation $I$) equals water out (evapotranspiration $\ET$, drainage $D$ below 1\,m, runoff $Q$) plus the change of water stored, $\Delta S$: $P+I=\ET+D+\Delta S+Q$.
\emph{Example.} Coarse loam without irrigation: $P=464.4$, $\ET=466.8$, $D=6.1$\,mm; $\ET+D-P=8.5$\,mm came from storage.

\subsection{2. the storage release}
$\bar R=\ET+D-P-I$ is the water the soil gave away from its store. Starting the run wetter makes $\bar R$ larger; the experiment E9 does this on purpose with ten starting suctions.

\subsection{3. the memory fate}
Plot drainage against $\bar R$ for the runs with the same rainfall: the slope is the share of each released millimetre that drained, $\phi_D$. The same slope for transpiration, soil and other evaporation gives $\phi_T$, $\phi_{E_s}$, $\phi_{\mathrm{oth}}$, and the four add to 1 because the released water must leave by one of them.
\emph{Example (three runs).} If $\bar R=10$, 20, 30\,mm and $D=1$, 4, 7\,mm, the slope is $(7-1)/(30-10)=0.3$: 30\,\% of released water drains.

\subsection{4. why drainage wins in coarse soils, the rate law}
Released water leaves by two routes that compete: drainage, fast when the saturated conductivity $K_s$ is large, and withdrawal by roots and the surface, with a fixed rate $q^\ast$. The share that drains is $\phi_D=K_s/(K_s+q^\ast)$, with $q^\ast=71.5$\,cm\,d$^{-1}$.
\emph{Example.} Sand, $K_s=712.8$: $712.8/(712.8+71.5)=0.91$ (observed 0.85). Silty clay loam, $K_s=1.68$: $1.68/73.18=0.023$ (observed 0.022).

\subsection{5. the recoverable share of an evaporation saving}
When a mulch or a lower evaporation demand saves evaporation, the saved water goes to transpiration, to less irrigation, or to drainage. Only the part that does not drain is recovered: $\rho_R=1-\psi_D$.
\emph{Example.} Loamy sand, $\psi_D=0.80$ in the edge-case experiment: $\rho_R=0.20$; only a fifth of a saving is recovered.

\subsection{6. checking the model's own bookkeeping}
If the outputs of the model obey the balance of Part 1, the columns $P$, $I$, $\ET$ and $D$ of all runs are linked by one linear relation. The singular-value decomposition finds the combination of columns that is closest to zero; it gave coefficients within 5\,\% of $(1,1,-1,-1)$, so the method finds the balance without being told it.

\subsection{7. the unattributed evaporation}
$\ET$ minus transpiration, soil evaporation and canopy evaporation should be zero; in the run table it averaged 164\,mm a year (30\,\% of $\ET$). The full outputs show it is evaporation of ponded water, litter and snow, which the run table did not report separately.

\subsection{8. reproduce everything}
Sheet \texttt{EQ\_P1} computes Part 4, Part 5 and the storage release of each soil from the runs; the scripts of \nolinkurl{Manuscripts/reproducibility} regenerate every number of the article.

